\section{Lista zadań numer 4.}

\begin{problem}[1]

	Udowodnić, że jeśli rząd grupy jest liczbą pierwszą, to G jest cykliczna.
\end{problem}
\begin{solution}
	
	Ustalmy grupę G i załóżmy, że jej rząd n jest liczbą pierwszą. To znaczy \[
	\text{rząd}\left( G \right) = n\text{ oraz } n \text{ jest liczbą pierwszą} \ge 2
	.\] 

	Ustalmy dowolny element $g \in G$, różny od $e$. Takie g istnieje, bowiem rząd G $\ge 2$. Niech k jest rzędem elementu g. k jest skończone i $\le n$, bo jeśli by było inaczej, to mielibyśmy sprzeczność z tym, że rząd G to n.
	
	Wniosek z twierdzenia Lagrangea mówi nam. że k jest dzielnikiem n. Ponieważ n jest liczbą pierwszą więc jej dzielnikami są 1 i n. Jedynym elementem w grupie G, który ma rząd 1 jest element neutralny e, ale k jest różne od e, zatem k jako dzielnik n jest równe n. Przeto $\langle g \rangle $ jest podgrupą cykliczną grupy G o rzędzie n, czyli $\langle g \rangle = G$, a co za tym idzie, G jest cykliczna.
\hfill\qed
\end{solution}

\begin{problem}[2]
	Niech N będzie podgrupą grupy G. Udowodnić, że następujące warunki są równoważne.

	\begin{enumerate}[label=\alph*)]
		\item $N\unlhd G$ czyli dla każdego $g\in G, g^{-1}Ng \subseteq N$.
		\item Dla każdego $g \in G$, $gNg^{-1}\subseteq N$.
		\item  Dla każdego $g \in G$, $gNg^{-1}= N$.
		\item Dla każdego $g \in G$, $g^{-1}Ng = N$. 
		\item  Dla każdego $g \in G$, $gN = Ng$.  
	\end{enumerate}
\end{problem}
\begin{solution}

	Pokażę ciąg wynikań, mianowicie
	\[
	a)\implies b) \implies c) \implies d) \implies e) \implies a)
	.\] 
	

	\begin{itemize}
		\item $a) \implies b)$ 

			Załóżmy a) $N\unlhd G$ czyli dla każdego $g'\in G, g'^{-1}Ng' \subseteq N$.

			Ustalmy dowolne $g\in G$. Korzystając z a), weźmy $g' = g^{-1}$. Wtedy $N \supseteq g'^{-1}Ng' = \left( g^{-1} \right)^{-1}Ng^{-1} = gNg^{-1} = gNg^{-1}$

		\item $b) \implies c)$ 

			Załóżmy b), że dla każdego $g' \in G$, $g'Ng'^{-1}\subseteq N$
	
			Ustalmy dowolne $g\in G$ oraz $n\in N$.
			Z założenia b), dla $g' = g^{-1}$ $g^{-1}ng \in N$. Zatem $n = \left( e \right) n \left( e \right) = g\left( g^{-1} n g \right)g^{-1} \in gNg^{-1} $
                        
			Mamy podwójne zawieranie, więc $gNg^{-1}=N.$
		\item $c) \implies d)$

			Załóżmy c), że dla każdego $g' \in G$, $g'Ng'^{-1}= N$. 
			Ustalmy dowolne $g\in G$. Niech $g' = g^{-1}$

			
			Wtedy z c) mamy $g^{-1}Ng = g^{-1}N\left( g^{-1} \right)^{-1} = g'Ng'^{-1} = N $
		\item $d) \implies e)$

			Załóżmy d), że dla każdego $g'\in G$, $g'Ng'^{-1} = N$. 
			Ustalmy dowolne $g\in G$.
			Mamy następujące równoważności:

			\begin{align*}
				h \in gN & \iff \left(\text{z d) } N = g^{-1}Ng\text{, dla } g' = g^{-1}\right) \\
				h \in g\left( g^{-1}Ng \right) & \iff \\
				h \in \left( gg^{-1} \right) N g & \iff\\
				h \in Ng
			\end{align*}

			Zatem $gN = Ng$
		\item $e) \implies a)$.

			Załóżmy e), że dla każdego  $g' \in G$ $gN = Ng$

			Ustalmy dowolne $g\in G$ oraz $h\in g^{-1}Ng$.

			Z założenia e) mamy
			\[
			N = \left( g^{-1}g \right)N = g^{-1}\left( gN \right) = g^{-1}\left( Ng \right) = g^{-1}Ng \subseteq N
			.\] 
			Zatem $h \in N.$
	\end{itemize}

\end{solution}

\begin{problem}[3]

	Niech $f: G \rightarrow H$ będzie homomorfizmem grup. Udowodnić, że $\ker(f) \le G$ oraz $\im(f) \le H$.
\end{problem}
\begin{solution}

	Załóżmy, że $f:G\rightarrow H$ jest homomorfizmem grup.
	\begin{itemize}
		\item $\ker(f) \le  G$ 
			
			Z definicji 
			\[
			\ker(f) = \{g \in G: fg = e_{H}\} 
			.\] 
			Zauważmy, że
			\begin{enumerate}
				\item $\ker(f) $ jest zamknięty na działanie:

					Weźmy dowolne dwa elementy $g,h \in \ker(f)$. Pokażę, że $gh\in \ker(f)$. Zauważmy, że
					\[
						g,h \in \ker(f) \iff fg = e_H \text{ oraz } fh = e_H 
					.\] 					Zatem
					\[
					e_H = \left( e_H \right) \left( e_H \right) = \left( fg \right) \left( fh \right) = f\left( gh \right)  					.\] 
					Czyli $f\left( gh \right) = e_H$, a co za tym idzie $gh \in \ker(f) $
		\item $e_{G} \in \ker(f) $

			Niech $h = fe_G \in H$.
			Wtedy
			 \[
			h = fe_G = f\left( e_Ge_G \right) = f\left( e_G \right) f\left( e_G \right) = hh  
			.\]
			Zatem
			\[
			e_H = h^{-1}h = h^{-1}\left( hh \right) = \left( h^{-1}h \right) h = e_Hh = h
			.\]  czyli $e_H = h = fe_G$.

			Zatem  $e_G\in \ker(f) $.
			\end{enumerate}
		\item $\im(f) \le H$

			Z definicji
			\[
			\im(f)  = \{h \in H: \exists_{g\in G} \left( fg = h \right)  \} 
			.\] 

			Zauważmy, że
			\begin{enumerate}
				\item $\im(f)$ jest zamknięte na działanie.

					Weźmy dowolne $h_1,h_1 \in \im(f) $.
					Wtedy istnieją $g_1,g_2$, takie że $fg_1 = h_1, fg_2 = h_2$
					
					Zatem

					\[
					h_1g_2 = \left( fg_1 \right) \left( fg_2 \right) = f\left( g_1g_2 \right) \in \im(f) 
					.\] 
					Czyli $h_1h_2 \in \im(f) $
				\item $e_H \in \im(f)$ 
					
					Niech $h = fe_G$. Wtedy
					 \[
					h = fe_G = f\left( e_Ge_G \right) = \left( fe_G \right) \left( fe_G \right) = hh
					.\] 
					Zatem 
					\[
					e_H = h^{-1}h = h^{-1}\left( hh \right) = \left( h^{-1}h \right) h = e_Hh = h = fe_G \in \im(f)  
					.\] 
			\end{enumerate}
	\end{itemize}
\end{solution}
\begin{problem}[4]

	Udowodnić, że jeśli $H\le G$ oraz $[G:H] = 2$, to $H\unlhd G$. 
\end{problem}
\begin{solution}

	Zauważmy, że warstwy H dzielą grupę G na rozłączne podzbiory- tworzą partycję zbioru G. Indeks H w G wynosi 2, a więc G jest podzielone na dwie rozłączne części. 

	Ponieważ $H\le G$, więc jedną z warstw jest samo $H = eH = He$.
        Drugą warstwą zatem musi być dopełnienie H w G oznaczmy je $H^{c}$. 

	$H\unlhd G$ jest równoważne temu, że dla dowolnego $g \in G$, $gH = Hg$. Ustalmy więc dowolne  $g\in G$ i rozważmy dwa przypadki:
	\begin{enumerate}
		\item $g\in H$

			Ponieważ $H\le G$ więc H jest zamknięte na działanie, czyli $gH = H = Hg$. Zatem  $H\unlhd G$.
		\item $g \not\in H$ 

			Ponieważ $H\le G$ więc $gH$ nie może być równe  $H$. Analogicznie $Hg \neq H$. Ponieważ, jedynym zbiorem został $H^{c}$, więc $gH = H^{c} = Hg$
	\end{enumerate}


\begin{proof}

	Aby udowodnić, że $H$ jest podgrupą normalną w $G$, musimy pokazać, że dla każdego elementu $g\in G$ zachodzi równość warstwy lewostronnej i prawostronnej: $gH = Hg$

Z założenia $[G:H] = 2$ wiemy, że istnieją dokładnie dwie różne lewostronne warstwy podgrupy H w G. Z lematu o warstwach wiemy, że warstwy te tworzą partycję zbioru G.

Jedną z tych warstw jest zawsze sama podgrupa $H$ (gdyż $H = eH$ dla $e \in H$). Ponieważ warstwy są rozłączne i pokrywają całą grupę, druga warstwa musi być zbiorem wszystkich pozostałych elementów, czyli dopełnieniem $G \setminus H$. Zatem zbiór lewych warstw to $\{H, G \setminus H\}$.

Analogiczne rozumowanie stosuje się do warstw prawostronnych. Istnieją dokładnie dwie prawe warstwy, którymi muszą być $H$ oraz $G \setminus H$.

Rozważmy teraz dwa możliwe przypadki dla dowolnego $g \in G$:

\begin{enumerate}
    \item \textbf{Przypadek 1: $g \in H$}
    
        Wtedy, z własności podgrupy (zamkniętość na działanie), warstwa lewostronna $gH$ jest równa $H$. Podobnie, warstwa prawostronna $Hg$ jest równa $H$.
        Zatem w tym przypadku zachodzi równość:
        \[
            gH = H = Hg.
        \]

    \item \textbf{Przypadek 2: $g \notin H$}
    
        Ponieważ $g$ nie należy do podgrupy $H$, warstwa lewostronna $gH$ nie może być równa $H$ (gdyby $gH=H$, to $g = ge \in gH = H$, co jest sprzeczne z naszym założeniem). Skoro istnieją tylko dwie lewe warstwy, to $gH$ musi być tą drugą, czyli:
        \[
            gH = G \setminus H.
        \]
        Analogicznie, warstwa prawostronna $Hg$ nie może być równa $H$. Zatem musi być równa $G \setminus H$:
        \[
            Hg = G \setminus H.
        \]
        Porównując oba wyniki, otrzymujemy równość:
        \[
            gH = G \setminus H = Hg.
        \]
\end{enumerate}

Ponieważ równość $gH = Hg$ zachodzi dla każdego elementu $g \in G$ (zarówno dla $g \in H$, jak i $g \notin H$), podgrupa $H$ spełnia definicję podgrupy normalnej.
\end{proof}

\end{solution}
\begin{problem}[5]

	Udowodnić, że $T_{n}\left( \R \right) $ nie jest dzielnikiem normalnym w $GL_{n}\left( \R \right) $.
\end{problem}

	Najpierw oznaczenia:
	\begin{itemize}
		\item $T_{n}\left( \R \right)$ jest to grupa macierzy dolnotrójkątnych o wyrazach rzeczywistych i rozmiarze $n\times n$. $T$ pochodzi od angielskiego \textbf{Triangular} 

			Analogicznie mamy $UT_{n}\left( \R \right)$ gdzie $UT$ pochodzi od angielskiego \textbf{Upper Triangular} 
		\item $GL_{n}\left( \R \right) $ jest to grupa macierzy odwracalnych o wyrazach rzeczywistych i rozmiarze $n\times n$. $GL$ pochodzi z angielskiego  \textbf{General Linear} 
	\end{itemize} 

\begin{solution}
	
	Aby pokazać, że $T_{n}\left( \R \right) \text{ nie jest dzielnikiem } GL_{n}\left( \R \right)$ musimy dla ustalonego $n$ wskazać takie $g\in GL_{n}\left( \R \right)$, że $gT_{n}\left( \R \right)g^{-1} \neq T_{n}\left( \R \right) $

	Zatem niech $n = 2$ oraz 
	
	\[t = \begin{pmatrix} 1 & 1\\0 & 1\end{pmatrix} \in T_{2}	\left( \R \right), g = \begin{pmatrix} 1 & 0\\ 1 & 0 \end{pmatrix},g^{-1} = \begin{pmatrix} 1 & -1 \\ 0 & 1 \end{pmatrix} \in GL_{n}\left( \R \right) \] 

	Wtedy $gtg^{-1} \in gT_{2}\left( \R \right)g^{-1}$ ale

		\[gt = \begin{pmatrix} 1 & 1 \\ 1 & 2 \end{pmatrix} 
		\text{ oraz } 
		\left( gt \right)g^{-1} = \begin{pmatrix} 0 & 1 \\ -1 & 2 \end{pmatrix}.\]

		Zatem \[gtg^{-1} \not\in T_{2}\left( \R \right)\] 
		czyli 
		\[gT_{2}\left( \R \right) g^{-1} \neq T_{2}\left( \R \right) .\]


\end{solution}

\begin{problem}[6]
	
	Korzystając z zasadniczego twierdzenia o homomorfizmie grup, udowodnić, że
	\begin{enumerate}[label=\alph*)]
		\item $\left( \R*, \cdot  \right)/\left( \{-1,1\}  \right) \cong \left( \R_{>0}, \cdot  \right)$,
		\item $\left( \C, + \right) /\Z \cong \left( \C*, \cdot  \right)$,
		\item $\left( \C, + \right) /\langle \exp^{\frac{2\pi i}{n}} \rangle  \cong \left( \C*, \cdot  \right)$,
	\end{enumerate}
	
\end{problem}
\begin{solution}

	\begin{enumerate}[label=\alph*)]
		\item Niech  $\varphi \left( x \right) = |x|$, czyli wartość bezwzględna.
			\begin{center}
		\begin{tikzcd}[row sep= 30pt, column sep=30pt]
			\left( \R^{*},\cdot  \right)  \arrow[dr, "\pi"] \arrow[rr, "\varphi"] & & \left( \R_{>0},\cdot  \right)   \\					     & \mfaktor{\left( \R^{*},\cdot  \right)}{\{-1,1\}}  \arrow[ur, "\psi", dashrightarrow] &  
		\end{tikzcd}
	\end{center}
	\item Niech $\varphi\left( z \right) = e^{2\pi i z}$
			\begin{center}
		\begin{tikzcd}[row sep= 30pt, column sep=30pt]
			\left( \C,+ \right)  \arrow[dr, "\pi"] \arrow[rr, "\varphi"] & & \left( \C^{*},\cdot  \right)   \\					     & \mfaktor{\left( \C, +  \right)}{\Z}  \arrow[ur, "\psi", dashrightarrow] &  
		\end{tikzcd}
	\end{center}
	\item Mamy taki diagram
			\begin{center}
		\begin{tikzcd}[row sep= 30pt, column sep=30pt]
			\left( \C^{*},\cdot  \right)  \arrow[dr, "\pi"] \arrow[rr, "\varphi"] & & \left( \C^{*},\cdot  \right)   \\					     & \mfaktor{\left( \C^{*},\cdot  \right)}{\langle \exp^{\frac{2\pi i}{n}} \rangle }  \arrow[ur, "\psi", dashrightarrow] &  
		\end{tikzcd}
	\end{center}



	\end{enumerate}
\end{solution}
\begin{problem}[7]
	
	cosniecos
\end{problem} 
\begin{solution}

Mamy $\Z_{6}$.
Szukamy jej wszystkich podgrup normalnych i podgrupy.
$Z_{6}$ jest abelowa więc wszystkie są normalne.

 \begin{itemize}
	\item $P_{1} = \{0\} \text{ rząd } = 1$ 
	\item $P_2 = \{0, 2, 4\} \text{ rząd }= 3$ 
	\item $P_{3} = \{0, 3\} \text{ rząd }= 2$ 
	\item $P_5 = \{0, 1, 5\} \text{ rząd }= 6$
\end{itemize}


	

\end{solution}





